\documentclass[11pt]{article}

\usepackage[margin=1in]{geometry}
\usepackage{amsmath,amssymb}
\usepackage{array}
\usepackage{booktabs}
\usepackage{caption}
\usepackage{hyperref}
\emergencystretch=2em

% T80-REOPEN: Title
% T78-SCOPE: title-hierarchy-demotion; source=T77-manifest,T70,T72,T77
\title{A Dual-Loop Teacher-Anchored Residual Calibration Framework for Drift-Adaptive GKP Decoding}
\author{}
\date{}

\begin{document}
\maketitle

% T80-REOPEN: Abstract
% T78-SCOPE: evidence-alignment; source=T77-manifest,T70,T72,T77
\begin{abstract}
% T77-SOURCE: T74-TBL-01; T74-TBL-02; T74-TBL-03; T74-TBL-07; review=T70,T76
Approximate Gottesman--Kitaev--Preskill (GKP) quantum error correction uses
continuous modular syndromes to infer displacement errors, but finite-energy
states and drifting operating conditions make a fixed correction law brittle.
We study a dual-loop decoding framework that preserves a deterministic affine
fast path, \(\Delta_t=K_t s_t+b_t\), while using recent syndrome histograms to
update the runtime parameters at a slower cadence.  The design keeps the
per-shot computation compatible with quantized fixed-point execution and makes
the adaptation target explicit: learning or calibration acts only on a
low-dimensional parameter surface, not on the full decoding decision.

A classical teacher estimator provides an interpretable anchor for the affine
parameters, and a lightweight CNN predicts only a bounded residual calibration
term, currently focused on the bias \(b_t\).  This differs from end-to-end
neural decoders and from outer-code soft-information decoders: the learned
module updates a physical-layer calibration surface, whereas the per-shot path
remains a deterministic affine rule.  The intended systems advantage is to
reduce logical failures under drift while keeping the fast boundary close to a
small fixed-point controller.

Under the locked four-scenario mock-backed software-HIL protocol, the
teacher-anchored hybrid residual branch is the winner in all four frozen
scenarios, with UKF as the runner-up throughout.  Bounded ablation and
multi-seed materials support only a descriptive reading of this result:
histogram-delta features matter, the instability pattern repeats across seeds,
and the tested lower-clip intervention is mixed and mostly harmful, so the
current mechanism story remains non-causal.  A separately labeled
supplement-side statistical-calibration extension lane further shows that a
simple non-neural estimator can be strong under the same affine runtime
contract, but its accepted closure remains no-promotion and without a unique
clean threshold.  The strongest current claim is therefore not that one CNN
architecture closes the decoding problem, but that drift-adaptive affine
calibration is a credible and hardware-bounded organizing principle for GKP
fast-path decoding.
\end{abstract}

\tableofcontents

% T80-REOPEN: Introduction
% T78-SCOPE: evidence-alignment; source=T77-manifest,T70,T72,T77
\section{Introduction}

Bosonic codes protect quantum information by encoding a logical qubit into the
Hilbert space of an oscillator.  The GKP code is a central example: ideal code
states form a periodic lattice in phase space, so small shift errors can be
diagnosed through modular syndrome measurements and corrected by displacement
operations \cite{gkp2001,grimsmo2021}.  Unlike binary stabilizer syndromes, GKP
syndromes are analog.  Their continuous values carry information not only
about the likely displacement, but also about how close the correction is to a
lattice decision boundary.

The practical difficulty is that physical GKP states are approximate
finite-energy states.  Their comb peaks have finite width and an envelope, and
syndrome measurements are further affected by measurement inefficiency, noisy
auxiliary states, circuit-level imperfections, oscillator loss, and calibration
error \cite{grimsmo2021,roy2025}.  As a result, the effective noise model seen
by the decoder can drift over time.  A correction law tuned for one operating
point may become systematically mismatched after bias, variance, or covariance
changes.

Recent GKP and bosonic-decoding work makes clear that analog soft information
is valuable.  Surface-GKP and concatenated bosonic-QLDPC decoders use
continuous syndrome information to set matching weights, belief-propagation
messages, or effective outer-code priors \cite{noh2020,noh2022,raveendran2022,
berent2024,borah2025}.  In parallel, hardware-facing QEC studies show that
learned modules are most useful when their latency role and integration
boundary are explicit \cite{chamberland2026,stein2026}, and calibration work
emphasizes that decoder priors must remain aligned with measured device
statistics \cite{spitz2018,wagner2021,dgr2023,chen2022,sivak2024}.  These two
lessons motivate the present formulation.

Our starting point is that drift adaptation should enter the GKP layer through
a bounded runtime contract rather than a free-form per-shot decoder.  The fast
loop is intentionally simple:
\begin{equation}
    \Delta_t = K_t s_t + b_t ,
\end{equation}
where \(s_t\) is the measured two-quadrature GKP syndrome and
\((K_t,b_t)\) are quantized runtime parameters.  A slower loop observes recent
syndrome histograms, estimates an effective noise state, and stages updated
parameters into the runtime bank.  The online correction path therefore remains
a small affine operation, while nonstationary estimation is handled
asynchronously through teacher-guided or calibration-driven updates.

This separation gives the method a specific systems role.  Compared with a
fixed affine decoder, the parameters can track drift.  Compared with a full
neural decoder, the per-shot computation remains smaller, more interpretable,
and easier to validate in fixed-point form.  Compared with outer-code
soft-information decoders, the adaptation target is the physical-layer GKP
displacement rule itself rather than a matching graph or LDPC message schedule.
The learned branch is therefore not the real-time decoder; it is one possible
slow-loop calibration module operating on a low-dimensional affine parameter
surface.

This also positions the work relative to FPGA-QEC practice.  Much of the
real-time hardware literature focuses on decoding binary stabilizer syndromes
for surface codes, repetition-code experiments, or quantum-LDPC memories using
lookup tables, union-find variants, matching or greedy approximations,
clustering, neural decoders, or message-passing hardware
\cite{lilliput2022,afs2022,helios2023,astrea2023,collision2025,
caune2024,yang2026,maurer2025}.  Those results establish the standard for
closed-loop latency and resource reporting, but they do not directly address a
GKP physical correction layer whose measured syndrome is analog, modular, and
shaped by finite-energy state preparation.  The present architecture is
therefore complementary: it provides a low-dimensional calibration layer that
could sit before, inside, or alongside a larger stabilizer or LDPC decoding
stack while keeping the fast boundary as a small fixed-point affine
displacement.

The current note keeps the evidence hierarchy explicit.  The locked frozen
software-HIL benchmark is the main result layer.  Feature and teacher
ablations, six-seed mechanism materials, and the separately labeled
statistical-calibration extension lane are supporting layers that help explain
or contextualize the benchmark but do not replace it.  Training/material
regeneration, isolated true \texttt{.tflite} execution, and real-board gate or
provenance artifacts remain layered boundary evidence rather than coequal
performance results.

% T78-SCOPE: evidence-alignment; source=T77-manifest,T70,T72,T77
\section{Summary of Contributions}

The work can be organized around five contributions.

\paragraph{1. A two-timescale adaptive affine formulation for approximate GKP
decoding.}
The per-shot correction path is formulated as a bounded affine estimator
\(\Delta_t=K_t s_t+b_t\).  The expensive adaptation problem is moved into a
slow loop that estimates the current effective noise state from recent syndrome
statistics.  This links the continuous-syndrome GKP picture to a runtime
contract that can be implemented as a quantized matrix-vector operation.

\paragraph{2. A teacher-anchored residual calibration strategy.}
A classical teacher estimator produces a stable baseline
\((K_t^{\rm teacher},b_t^{\rm teacher})\).  The CNN predicts a small residual,
\begin{equation}
    \delta b_t=f_\phi(H_{t-c+1:t},\Delta H_t,z_t^{\rm teacher}),
\end{equation}
and the committed parameters are
\begin{equation}
    K_t=K_t^{\rm teacher}, \qquad
    b_t={\rm EMA}\!\left(b_t^{\rm teacher}+\delta b_t\right).
\end{equation}
Learning is therefore used to calibrate a low-dimensional control surface, not
to replace the GKP decoder end-to-end.

\paragraph{3. A deployment-aware runtime architecture.}
The architecture exposes fixed-point quantization, clipping, saturation
diagnostics, parameter smoothing, stale-parameter behavior, and double-buffered
stage-and-commit updates.  These constraints become measurable objects in the
benchmark rather than assumptions added after the algorithm is designed.

\paragraph{4. A benchmark suite for adaptive affine calibration.}
% T77-SOURCE: T74-TBL-01; T74-TBL-02; T74-TBL-03; T74-TBL-07; review=T70,T76
The mainline numerical study compares filtering baselines and learned residual
variants under four frozen drift scenarios.  Bounded ablation and multi-seed
follow-ups then support a conservative reading of the gains.  A separately
labeled supplement-side statistical calibration extension lane asks whether a simple
non-neural rule can also be strong inside the same runtime contract, but it
remains a supplement-side no-promotion result rather than a promoted mainline
comparator.

\paragraph{5. A bridge between GKP analog calibration and real-time decoder
hardware practice.}
The formulation exposes a clear interface between continuous GKP syndrome
statistics and a hardware-style decoder contract: the slow side may use
histograms, teachers, learned residuals, or statistical estimators, but the
fast side only consumes bounded quantized parameters.  This makes the proposed
method more directly comparable with FPGA-QEC systems than an unconstrained
offline neural decoder, while targeting a GKP-specific problem that binary
stabilizer decoders do not solve on their own.

\subsection{Metric-level advantages}

For reviewer-facing comparison, the advantages should be stated by metric
rather than by architecture alone.  Table~\ref{tab:metric-advantages}
separates the evidence currently supported by software-HIL experiments from
the hardware-facing advantages that follow from the fast-path contract but
still require board-level measurement.

\begin{table}[ht]
\centering
\scriptsize
\renewcommand{\arraystretch}{1.06}
\begin{tabular}{@{}>{\raggedright\arraybackslash}p{0.18\linewidth}>{\raggedright\arraybackslash}p{0.24\linewidth}>{\raggedright\arraybackslash}p{0.29\linewidth}>{\raggedright\arraybackslash}p{0.20\linewidth}@{}}
\toprule
Metric & Advantage & Architectural reason & Evidence boundary \\
\midrule
Logical-error proxy &
Lower error under controlled drift than filtering baselines; statistical
calibration gives an even stronger non-neural reference point. &
Both learned and statistical branches update the same physical-layer affine
surface, so adaptation acts directly on the displacement rule used by the GKP
fast loop. &
Supported in the four-scenario software-HIL suite; not a paper-grade expanded
benchmark. \\
\addlinespace
Per-shot latency &
Low latency by design on the correction path. &
The fast loop evaluates only a \(2\times2\) affine map, bias addition,
clipping, quantization, and parameter-bank selection; CNN or statistical
estimation is kept off the per-shot critical path. &
Architectural claim; board timing and worst-case latency remain a hardware
measurement target. \\
\addlinespace
FPGA resource and cost &
Expected small fast-path resource footprint and lower verification burden than
placing a neural or graph decoder in the real-time path. &
The real-time datapath needs only fixed-point multiply-adds, saturation logic,
and a compact parameter bank; the slow estimator can run less frequently on a
host, embedded processor, or offline update service. &
Resource use and energy/cost are not yet measured on a real device. \\
\addlinespace
Engineering implementation &
The control boundary is simple to integrate, test, and fail safely. &
Quantized \((K,b)\), clipping limits, stale-parameter handling, saturation
counters, and double-buffered stage-and-commit updates form explicit runtime
contracts. &
Software-HIL validates the interface shape; real-board execution is outside
the current evidence layer. \\
\addlinespace
Drift robustness &
Adaptation can track mean, variance, and orientation changes in the effective
GKP noise state. &
Histogram windows and teacher/statistical summaries update the affine
parameters over time instead of freezing one calibration point. &
Supported for the current drift suite; broader drift families and holdout
protocols remain benchmark-expansion targets. \\
\addlinespace
Modularity &
The best slow-loop estimator can change without changing the fast path. &
Teacher residual CNN, statistical calibration, temporal models, and
gain-scheduled parameter banks all map to the same committed \((K,b)\)
interface. &
This is the reason a strong statistical-calibration result strengthens the
project-level contract rather than invalidating it. \\
\addlinespace
Compatibility with larger QEC stacks &
The method can serve as a GKP physical-layer calibration module before an
outer stabilizer, surface-GKP, or LDPC-GKP decoder. &
It improves the physical displacement correction and residual analog
statistics instead of replacing the outer-code decoder decision rule. &
System-level concatenated-code gains remain future-work evidence. \\
\bottomrule
\end{tabular}
\caption{Metric-level advantage summary.  The strongest current evidence is
logical-error reduction in software-HIL.  Latency, resource, and cost
advantages are architecture-level consequences of the affine fast path and
should be measured on hardware before being written as deployment results.}
\label{tab:metric-advantages}
\end{table}

\section{Brief Review of the GKP Code}

\subsection{Ideal and approximate code states}

The square-lattice GKP code encodes a qubit into an oscillator by using
periodic structure in the \(q\) and \(p\) quadratures \cite{gkp2001}.  In the
convention used here, the lattice constant is
\begin{equation}
    \lambda=\sqrt{2\pi}.
\end{equation}
An ideal logical state may be represented heuristically as an infinite comb,
for example
\begin{equation}
    |\bar 0\rangle \propto \sum_{n\in\mathbb{Z}} |n\lambda\rangle_q .
\end{equation}
Ideal comb states are not physical because they require infinite energy.
Approximate GKP states replace the infinitely sharp comb with finitely squeezed
peaks and an envelope.  Finite-energy structure is central to practical GKP
QEC rather than a minor implementation detail \cite{grimsmo2021}: it
determines the intrinsic syndrome uncertainty and contributes an effective
displacement-like noise.

\subsection{Syndrome measurement as modular displacement information}

Let the accumulated phase-space displacement error before correction be
\begin{equation}
    e_t=\begin{bmatrix}e_{q,t}\\e_{p,t}\end{bmatrix}.
\end{equation}
An ideal GKP syndrome observes this displacement modulo the lattice:
\begin{equation}
    s_t = e_t \bmod \lambda,
    \qquad
    s_{q,t},s_{p,t}\in[-\lambda/2,\lambda/2).
\end{equation}
In a finite-energy and noisy-measurement setting, the measured syndrome is
more accurately described as
\begin{equation}
    \tilde{s}_t={\rm mod}(e_t,\lambda)+\eta_t^{\rm meas}+\eta_t^{\rm GKP}.
\end{equation}
Here \(\eta_t^{\rm meas}\) represents measurement and auxiliary-state
contributions, while \(\eta_t^{\rm GKP}\) summarizes finite-squeezing
uncertainty.  The decoder therefore observes a noisy modular representative,
not the absolute displacement.

\subsection{Local affine decoding and its limits}

The modulo structure makes exact GKP decoding nonlinear and branch dependent.
Near a single lattice branch, however, a Gaussian approximation gives a useful
local model.  If \(e\) and \(s\) are jointly Gaussian,
\begin{equation}
\begin{bmatrix}e\\s\end{bmatrix}
\sim
\mathcal{N}
\left(
\begin{bmatrix}\mu_e\\\mu_s\end{bmatrix},
\begin{bmatrix}
\Sigma_{ee} & \Sigma_{es}\\
\Sigma_{se} & \Sigma_{ss}
\end{bmatrix}
\right),
\end{equation}
then the linear-MMSE estimate is
\begin{equation}
    \hat e=\mu_e+\Sigma_{es}\Sigma_{ss}^{-1}(s-\mu_s)=Ks+b,
\end{equation}
where
\begin{equation}
    K=\Sigma_{es}\Sigma_{ss}^{-1},
    \qquad
    b=\mu_e-K\mu_s.
\end{equation}
This derivation explains why an affine fast path is a plausible low-latency
approximation.  It also marks the boundary of the approximation.  Near lattice
decision boundaries, the posterior can be multimodal; maximum-likelihood,
closest-lattice-point, or wrapped-Gaussian decoders can outperform a single
global affine rule \cite{fukui2018,roy2025}.

\subsection{Logical failure criterion}

After correction, the residual displacement is wrapped into the GKP
fundamental cell.  A logical error occurs when the residual crosses a logical
decision boundary, for example
\begin{equation}
    |r_{q,t}|>\lambda/2 \quad \Rightarrow \quad X_L\ {\rm error},
\end{equation}
and similarly for the \(p\)-quadrature.  Closed-loop logical error probability
is therefore the central metric; parameter-regression error is useful only when
it correlates with the downstream correction outcome.

\section{Noise and Drift Model}

\subsection{Effective noise state}

The slow loop does not attempt to infer every microscopic hardware parameter.
It uses a low-dimensional effective state,
\begin{equation}
    \theta_t^{\rm noise}
    =
    (\sigma_t,\mu_{q,t},\mu_{p,t},\vartheta_t),
\end{equation}
where \(\sigma_t\) controls the displacement scale, \(\mu_{q,t}\) and
\(\mu_{p,t}\) represent mean bias, and \(\vartheta_t\) captures covariance-axis
rotation.  This compact parameterization can be mapped to affine runtime
parameters and estimated from finite histogram windows.

\subsection{Noise sources represented by the model}

The model is an effective displacement-noise model with additional measurement
uncertainty.  It absorbs several physical effects:
\begin{enumerate}
    \item finite-energy GKP peak width and envelope effects;
    \item Gaussian random displacement noise in \(q\) and \(p\);
    \item biased displacement means caused by calibration offsets;
    \item anisotropic or rotated covariance caused by asymmetric noise;
    \item syndrome measurement noise and noisy auxiliary-state contributions.
\end{enumerate}
This is narrower than a full circuit-level GKP noise model, but it is aligned
with the affine fast-path contract.

\subsection{Drift scenarios}

The benchmark suite uses four drift scenarios:
\begin{enumerate}
    \item \texttt{static\_bias\_theta}: a biased and rotated operating point;
    \item \texttt{linear\_ramp}: slowly changing effective noise parameters;
    \item \texttt{step\_sigma\_theta}: a sudden variance/orientation change;
    \item \texttt{periodic\_drift}: periodic nonstationarity.
\end{enumerate}
These scenarios separate different adaptation demands: steady calibration,
tracking, shock response, and periodic following.

\section{Model Architecture}

\subsection{Fast loop: affine fixed-point decoder}

The fast loop receives the current measured syndrome \(s_t\), reads the active
parameter bank, and computes
\begin{equation}
    s_t^{\rm clip}={\rm clip}(s_t,-s_{\max},s_{\max}),
\end{equation}
\begin{equation}
    \Delta_t^{\rm raw}=K_t s_t^{\rm clip}+b_t,
\end{equation}
\begin{equation}
    \Delta_t=Q\!\left({\rm clip}(\Delta_t^{\rm raw},
    -\Delta_{\max},\Delta_{\max})\right),
\end{equation}
where \(Q(\cdot)\) denotes fixed-point quantization.  The implementation uses
an FPGA-oriented Q4.20 representation for syndrome values and runtime
parameters.  The fast loop also records diagnostic counters, including
histogram-input saturation, correction saturation, overflow, aggressive
parameter events, commit counts, and fallback-related signals.

\subsection{Parameter mapping}

Given an estimated noise state, the parameter mapper constructs an error
covariance
\begin{equation}
    C =
    R(\vartheta)
    \begin{bmatrix}
        \sigma_q^2 & 0\\
        0 & \sigma_p^2
    \end{bmatrix}
    R(\vartheta)^\top ,
\end{equation}
and a measurement covariance
\begin{equation}
    R_{\rm meas}=(\sigma_{\rm meas}^2+\Delta_{\rm eff}^2)I.
\end{equation}
The raw gain is
\begin{equation}
    K_{\rm raw}=C(C+R_{\rm meas})^{-1},
\end{equation}
followed by gain clipping or scaling.  The bias target is
\begin{equation}
    b_{\rm target}=\alpha(I-K_{\rm target})\mu,
    \qquad
    \mu=\begin{bmatrix}\mu_q\\\mu_p\end{bmatrix}.
\end{equation}
Finally, exponential smoothing produces staged parameters:
\begin{equation}
    K_t=(1-\beta)K_{t-1}+\beta K_{\rm target},
    \qquad
    b_t=(1-\beta)b_{t-1}+\beta b_{\rm target}.
\end{equation}

\subsection{Teacher estimators}

The teacher family provides stable estimates of \(\theta_t^{\rm noise}\) from
recent syndrome history.  A simple teacher computes moments over a histogram
window.  Stronger teachers such as EKF, UKF, RLS, or particle-filter variants
add state-space assumptions over time.  Abstractly,
\begin{equation}
    \hat{\theta}_t^{\rm teacher}={\rm Teacher}(H_{1:t}),
\end{equation}
and
\begin{equation}
    (K_t^{\rm teacher},b_t^{\rm teacher})
    =
    {\rm ParamMapper}(\hat{\theta}_t^{\rm teacher}).
\end{equation}
The teacher keeps the learned branch near an interpretable calibration surface
and creates meaningful ablation baselines.

\subsection{CNN residual branch}

The CNN consumes a short context of normalized syndrome histograms and
histogram deltas:
\begin{equation}
    X_t^{\rm hist}=[H_{t-c+1},\ldots,H_t,\Delta H_{t-c+2},\ldots,\Delta H_t].
\end{equation}
Teacher-side features may include
\begin{equation}
    z_t^{\rm teacher}
    =
    (\hat{\theta}_t^{\rm teacher},
     K_t^{\rm teacher},
     b_t^{\rm teacher},
     \Delta b_t^{\rm teacher},\ldots).
\end{equation}
The gated branch narrows this to a small set of teacher-\(b\) and
teacher-\(\Delta b\) scalars.  The learned branch predicts
\begin{equation}
    \widehat{\delta b}_t=f_\phi(X_t^{\rm hist},z_t^{\rm teacher}),
\end{equation}
then clips the result:
\begin{equation}
    \delta b_t={\rm clip}(s_b\widehat{\delta b}_t,-b_{\max},b_{\max}).
\end{equation}
The committed parameter is
\begin{equation}
    K_t=K_t^{\rm teacher},
    \qquad
    b_t={\rm EMA}(b_t^{\rm teacher}+\delta b_t).
\end{equation}

\subsection{Statistical calibration branch}

The same affine runtime contract also supports a non-neural statistical
calibration branch.  This branch estimates a bounded bias correction directly
from recent syndrome statistics and emits runtime parameters through the same
clipping, smoothing, and stage-and-commit path.  Its role is both practical and
scientific: it tests whether the benefit comes from a general histogram-driven
calibration principle rather than from the CNN architecture alone.

\subsection{Stage-and-commit runtime contract}

The slow loop does not directly mutate active fast-loop parameters.  It stages
a candidate \((K,b)\) into an inactive bank and commits it at a safe epoch
boundary:
\begin{equation}
(K_t,b_t)=
\begin{cases}
(K^A,b^A), & t<t_{\rm commit},\\
(K^B,b^B), & t\ge t_{\rm commit}.
\end{cases}
\end{equation}
This contract allows stale-parameter effects, update latency, commit success,
rollback/fallback behavior, and fixed-point stability to be measured
separately from logical performance.

% T80-REOPEN: Related Work / positioning
% T78-SCOPE: evidence-alignment; source=T77-manifest,T70,T72,T77
\section{Relationship to Existing Work}

\subsection{GKP and bosonic analog soft information}

GKP analog syndrome information is already known to be useful.  Analog-QEC,
surface-GKP, and concatenated bosonic-QLDPC decoders use continuous
measurement information to set decoder weights or soft messages
\cite{fukui2018,noh2020,noh2022,raveendran2022,berent2024,borah2025}.  The
distinction here is not the value of analog information itself, but the layer
at which it is consumed: recent syndrome statistics are compressed into
physical-layer affine GKP fast-path parameters rather than outer-code matching
weights or LDPC messages.

\subsection{Adaptive priors and syndrome-statistics estimation}

Adaptive-weight and syndrome-statistics methods estimate decoder weights or
noise parameters from measured data \cite{spitz2018,wagner2021,dgr2023}.
Calibrated decoders and decoder-prior optimization show that prior mismatch can
materially affect logical performance \cite{chen2022,sivak2024}.  This work
takes the same lesson seriously, but maps the adaptation target to the GKP
affine parameters \((K,b)\) and constrains the update through an explicit
stage-and-commit runtime contract.

\subsection{Learned low-latency QEC modules}

AI pre-decoders, high-accuracy neural decoders, and calibration-conditioned
FiLM decoders demonstrate a systems-oriented pattern: a learned module is most
useful when its input/output contract, latency role, and integration with a
classical decoder are explicit \cite{chamberland2026,stein2026,
google2023decoder}.  The learned module here is narrower.  It is not the
per-shot decoder, but a slow calibration component that updates
low-dimensional affine parameters while the fast boundary remains fixed-point
affine logic.

\subsection{Real-time and FPGA QEC decoders}

Real-time QEC decoders and FPGA-integrated systems have reported low-latency
hardware decoding in stabilizer-code settings \cite{caune2024,yang2026,
ziad2024,maurer2025}.  Hardware-oriented decoder designs also include
lookup-table approaches, union-find architectures, matching and greedy
approximations, clustering decoders, and FPGA-tailored quantum-LDPC
message-passing systems \cite{lilliput2022,afs2022,helios2023,astrea2023,
collision2025,qldpcfpga2025}.  Real-time bosonic QEC has also been
demonstrated in experimental feedback systems \cite{sivak2023,lachance2024}.
These works set the standard for full deployment claims: closed-loop timing,
worst-case latency, resource use, and hardware integration.

The contribution here is intentionally narrower and complementary.  It does
not aim to replace surface-code, repetition-code, or qLDPC hardware decoders.
Instead, it targets the GKP physical-layer correction problem that precedes
many concatenated-code decoding tasks: how should a device update the analog
displacement rule when the finite-energy GKP syndrome distribution drifts?  In
that setting, a fixed binary graph decoder is not the natural first
abstraction, and an unconstrained neural network on the per-shot path is not a
compelling hardware endpoint.  A staged affine parameter bank provides a
smaller and more testable contract: update \((K,b)\) slowly from recent analog
statistics, then execute only \(K s+b\) at the fast boundary.

\subsection{Advantages relative to existing QEC decoder approaches}

The proposed architecture is best viewed as a middle layer between physical
GKP syndrome acquisition and higher-level QEC decoding.  Its advantages over
nearby approaches are architectural and evidence-bounded rather than claims of
completed board-level deployment or promoted comparator closure.

\paragraph{Preserving analog GKP information at the right layer.}
Outer-code GKP methods already show that analog syndrome values improve
matching weights or belief-propagation messages.  This work uses the same
principle earlier in the stack: the analog histogram updates the physical GKP
displacement rule itself.  This can reduce prior mismatch before the
information is handed to a surface-code, LDPC, or other outer decoder.

\paragraph{Drift adaptation without a neural real-time critical path.}
End-to-end neural decoders can be expressive, but they make worst-case latency,
fixed-point validation, and fallback behavior harder to reason about.  Here,
learning is optional and narrow: a CNN residual branch or, in a separately
bounded extension lane, a statistical calibration branch updates a small
parameter surface, while the real-time path remains deterministic affine
logic.  This separation is expected to be useful when future hardware systems
need both adaptation and bounded per-shot timing.

\paragraph{A replaceable calibration module rather than a single decoder
commitment.}
The same runtime contract can host a teacher estimator, a learned residual
branch, or a non-neural statistical calibration rule.  The bounded
statistical-calibration extension lane therefore supports the design story
without changing the accepted hierarchy: the important object is the
histogram-driven affine calibration contract, not one particular CNN
architecture, and this note keeps that lane supplement-side rather than as a
promoted mainline comparator.

\paragraph{Explicit deployment diagnostics.}
The framework makes quantization, clipping, saturation, stale-parameter use,
commit counts, fallback behavior, and update latency visible benchmark
objects.  This is closer to FPGA-QEC practice than reporting only offline
regression error, and it provides a path to future resource and timing
measurements once a real hardware host is available.

\paragraph{Compatibility with larger QEC stacks.}
Because the fast path emits a physical displacement correction rather than a
complete outer-code decision, the method can in principle be combined with
surface-GKP, bosonic-QLDPC, or conventional stabilizer-code decoders.  Its
expected role is to keep the GKP layer calibrated under drift so that
downstream decoders receive better-conditioned residual information.

% T80-REOPEN: Experimental Setup
\section{Experimental Setup}

\subsection{Software-HIL protocol}

% T77-SOURCE: T74-TBL-01; review=T24,T25
The mainline numerical experiments use a locked mock-backed software-HIL
protocol that preserves the fast loop, slow loop, parameter-bank, and
stage-and-commit interfaces while executing the underlying experiment in
software.  The authoritative frozen benchmark is the `T24` protocol: four
drift scenarios, five mainline modes, paired seeds, and two repeats under one
fixed comparison contract.  This protocol is sufficient for comparing adaptive
calibration laws under matched drift trajectories and shared random structure,
but it does not measure board-level latency, hardware resource use, or
deployment closure.  Later ablation, multi-seed, and extension-lane materials
are therefore interpreted relative to this anchor rather than as a rewritten
benchmark.

\subsection{Scenarios and modes}

The four drift scenarios are:
\begin{center}
\small
\begin{tabular}{ll}
\toprule
Scenario key & Interpretation \\
\midrule
\texttt{static\_bias\_theta} & biased and rotated operating point \\
\texttt{linear\_ramp} & slowly changing effective noise parameters \\
\texttt{step\_sigma\_theta} & sudden variance/orientation change \\
\texttt{periodic\_drift} & periodic nonstationarity \\
\bottomrule
\end{tabular}
\end{center}

% T77-SOURCE: T74-TBL-01; T74-TBL-02; T74-TBL-03; T74-TBL-07; review=T70,T76
The mainline comparison includes five modes: EKF, UKF, a constant-mean affine
baseline, an RLS residual-\(b\) baseline, and the learned hybrid residual
branch.  Additional bounded materials then cover feature and teacher
ablations, six-seed descriptive mechanism snapshots, and a separately labeled
statistical-calibration extension lane under the same affine runtime contract.
The reported metric is \(\texttt{final\_ler\_mean}\), where lower values
indicate fewer logical failures under the fixed software-HIL protocol.

% T80-REOPEN: Numerical Results
\section{Numerical Results and Benchmark Plan}

\subsection{Four-scenario affine benchmark}

% T77-SOURCE: T74-TBL-01; T75-FIG-M01; T76-PREVIEW-M01; review=T76
Table~\ref{tab:five-mode-benchmark} records the locked frozen benchmark that
also underlies the paper-facing main-text visual summary.  Under the locked
`T24` protocol, the hybrid residual branch is the winner in all four frozen
scenarios, with UKF as the runner-up throughout.  This remains the
authoritative mainline result layer of the current note.  It supports a
bounded mock-backed software-HIL ranking only, and it should not be retold as
an expanded benchmark, a \texttt{.tflite} ranking, or a board-level ranking.

\begin{table}[ht]
\centering
\small
\begin{tabular}{lrrrrr}
\toprule
Scenario & EKF & UKF & Const.-\(\mu\) & RLS-\(b\) & Hybrid-\(b\) \\
\midrule
\texttt{static\_bias\_theta} & 0.838110 & 0.825370 & 0.836658 & 0.837577 & \textbf{0.810902} \\
\texttt{linear\_ramp} & 0.819200 & 0.811201 & 0.816911 & 0.819373 & \textbf{0.787755} \\
\texttt{step\_sigma\_theta} & 0.822365 & 0.811548 & 0.819784 & 0.821493 & \textbf{0.788800} \\
\texttt{periodic\_drift} & 0.832192 & 0.821558 & 0.829670 & 0.832334 & \textbf{0.806392} \\
\bottomrule
\end{tabular}
\caption{Four-scenario affine benchmark.  Entries are
\(\texttt{final\_ler\_mean}\); lower is better.}
\label{tab:five-mode-benchmark}
\end{table}

\subsection{Feature and teacher ablations}

% T77-SOURCE: T74-TBL-02; T74-TBL-03; T75-FIG-M02; T76-PREVIEW-M02; review=T57,T58,T76
Table~\ref{tab:feature-ablation} is best read as appendix-side support for the
frozen benchmark rather than as a replacement ranking table.  Removing
histogram deltas degrades performance relative to the full hybrid model,
indicating that temporal changes in the syndrome histogram matter inside this
fixed benchmark set.  Removing teacher-side channels also changes performance
materially, but the no-teacher-params row performing best means the current
evidence does not support a simple teacher-necessity claim.  The safest
reading is therefore descriptive: feature and teacher design matter, yet the
bounded ablation does not close the mechanism question.

\begin{table}[ht]
\centering
\small
\begin{tabular}{lrrr}
\toprule
Mode & Avg. LER & \(\Delta\) vs UKF & \(\Delta\) vs Hybrid Full \\
\midrule
\texttt{ukf} & 0.817382 & 0.000000 & +0.018837 \\
\texttt{hybrid\_full} & 0.798545 & -0.018837 & 0.000000 \\
\texttt{hybrid\_no\_hist\_deltas} & 0.826723 & +0.009341 & +0.028178 \\
\texttt{hybrid\_no\_teacher\_prediction} & 0.807251 & -0.010131 & +0.008706 \\
\texttt{hybrid\_no\_teacher\_params} & \textbf{0.749621} & -0.067761 & -0.048924 \\
\texttt{hybrid\_no\_teacher\_deltas} & 0.800329 & -0.017053 & +0.001784 \\
\bottomrule
\end{tabular}
\caption{Feature and teacher ablation summary.  Negative
\(\Delta\) means improvement over the reference.}
\label{tab:feature-ablation}
\end{table}

The next three lower-level note items preserve the accepted
statistical-calibration evidence as a supplement-side extension lane rather
than as a coequal Results pillar.

\subsubsection{Supplement-side statistical calibration extension lane}

% T77-SOURCE: T74-TBL-07; T74-SUP-01; review=T70
A separately labeled extension lane adds statistical calibration variants under
the same affine runtime contract.  Table~\ref{tab:statcalib-extension} records
the bounded comparison snapshot that motivated this lane: within that task-local
matrix, statistical calibration variants outperform UKF and the learned hybrid
branch.  However, the accepted closure deliberately does not promote this lane
into the mainline comparator set.  The locked five-mode frozen benchmark
remains the authoritative mainline result, and the statistical-calibration lane
is retained as supplement-side evidence for the broader adaptive-affine
contract.

\begin{table}[ht]
\centering
\small
\begin{tabular}{lrrr}
\toprule
Scenario & UKF & Hybrid-\(b\) & StatCalib \\
\midrule
\texttt{static\_bias\_theta} & 0.825370 & 0.810902 & \textbf{0.431708} \\
\texttt{linear\_ramp} & 0.811201 & 0.787755 & \textbf{0.467083} \\
\texttt{step\_sigma\_theta} & 0.811548 & 0.788800 & \textbf{0.460016} \\
\texttt{periodic\_drift} & 0.821558 & 0.806392 & \textbf{0.438751} \\
\bottomrule
\end{tabular}
\caption{Bounded statistical-calibration extension-lane comparison under the
same drift suite.  Entries are \(\texttt{final\_ler\_mean}\); lower is better.
This table does not by itself promote statistical calibration into the mainline
comparator set.}
\label{tab:statcalib-extension}
\end{table}

\subsubsection{Local sensitivity inside the statistical-calibration extension lane}

% T77-SOURCE: T74-TBL-07; T74-SUP-01; review=T66,T70
The local grid around the default setting shows that the extension-lane signal
is not tied to a single pathological point, but it also does not identify a
unique clean threshold choice.  The high-threshold variant has the best
aggregate mean LER by a negligible margin, whereas the default variant wins
three of four scenarios and has the best mean within-lane rank.  This
persistent tie is why the accepted closure remains no-promotion rather than
collapsing to one promoted statistical-calibration setting.

\begin{table}[ht]
\centering
\small
\begin{tabular}{lrrr}
\toprule
Variant & Mean LER & Mean Rank & Scenario Wins \\
\midrule
\texttt{statcalib\_high\_threshold} & \textbf{0.449241} & 1.75 & 1 \\
\texttt{statcalib\_default} & 0.449254 & \textbf{1.25} & \textbf{3} \\
\texttt{statcalib\_high\_scale} & 0.456477 & 3.00 & 0 \\
\texttt{statcalib\_low\_clip} & 0.481484 & 4.50 & 0 \\
\texttt{statcalib\_low\_scale} & 0.485129 & 4.50 & 0 \\
\bottomrule
\end{tabular}
\caption{Local sensitivity within the bounded statistical-calibration
extension lane.  Mean rank is computed within the five variants for each
scenario; the near-tie between default and high-threshold prevents a unique
promoted setting.}
\label{tab:statcalib-sensitivity}
\end{table}

\subsubsection{Per-scenario best variants and the no-promotion gate}

% T77-SOURCE: T74-TBL-07; T74-SUP-01; review=T69,T70
Table~\ref{tab:statcalib-scenario-best} reports the per-scenario winners inside
the extension lane.  These rows show why the lane remains scientifically
interesting, but they are not sufficient to override the no-promotion gate: the
default and high-threshold settings remain persistently tied across aggregate
and scenario-wise summaries, and the \texttt{static\_bias\_theta} winner also
carries a mixed-status caveat.  The safest use of these numbers is therefore
supplement-side documentation of the lane, not a claim that a mature promoted
comparator has been established.

\begin{table}[ht]
\centering
\small
\begin{tabular}{lrrr}
\toprule
Scenario & Best variant & Best LER & Gap vs Hybrid-\(b\) \\
\midrule
\texttt{static\_bias\_theta} & \texttt{high\_threshold} & 0.431684 & 0.379889 \\
\texttt{linear\_ramp} & \texttt{default} & 0.467094 & 0.321178 \\
\texttt{step\_sigma\_theta} & \texttt{default} & 0.458834 & 0.328883 \\
\texttt{periodic\_drift} & \texttt{default} & 0.439352 & 0.367232 \\
\bottomrule
\end{tabular}
\caption{Best statistical-calibration variant per scenario within the bounded
extension lane.  Gap is the hybrid residual LER minus the best statistical
calibration LER; this table does not remove the accepted no-promotion gate.}
\label{tab:statcalib-scenario-best}
\end{table}

\subsection{Mechanism probe for residual-b behavior}

% T77-SOURCE: T74-TBL-02; T74-TBL-03; T75-FIG-M02; T76-PREVIEW-M02; review=T58,T76
A paper-facing mechanism figure is available for the learned residual branch,
with appendix tables retaining the authoritative numeric snapshots.  The
six-seed evidence is descriptive rather than causal: the instability pattern
repeats broadly, while the tested lower-clip intervention is mixed and mostly
harmful.  This supports a cautious interpretation: residual amplitude
participates in some failure modes, but it is not a monotonic causal
explanation by itself, and the current evidence does not prove teacher
necessity.  Appendix tables provide bounded ablation, cross-seed numeric
snapshots, and provenance that support this conservative reading.

\subsection{Unseen drift generalization}

The current four scenarios cover steady bias, ramp tracking, step response, and
periodic drift.  A stronger benchmark would require unseen drift families such
as random-walk drift, \(1/f\)-like drift, burst measurement noise, coupled
bias-variance drift, and faster-than-window drift.  That expansion remains a
future benchmark lane rather than a completed claim in the present note; it is
precisely the kind of follow-up needed to distinguish genuine drift adaptation
from overfitting to a small set of hand-designed trajectories.

\subsection{Oracle and wrapped-Gaussian baselines}

The affine methods still need comparison against stronger theoretical
baselines: nearest-lattice hard decoding, static affine decoding, oracle
affine decoding with true noise state, and wrapped-Gaussian or
maximum-likelihood decoding in small settings.  Those baselines would separate
two kinds of loss that remain entangled here: the loss from constraining the
fast path to be affine, and the loss from estimating the noise state
imperfectly in the slow loop.

\subsection{Runtime, quantization, and fixed-point degradation}

% T77-SOURCE: T74-TBL-05; T74-SUP-03; review=T48
The deployment story depends on more than logical performance.  The current
paper-facing evidence therefore keeps runtime material at the boundary layer:
selected preserved \texttt{.tflite} artifacts have been executed in an isolated
current-host environment, but the broader fixed-point/runtime measurements
listed here remain future benchmark items.  Saturation and overflow rates,
commit latency, stale-parameter penalty, fallback frequency, and per-update
slow-loop cost are still needed to connect the mathematical affine rule to a
fully measured real-time decoder contract.

\subsection{Embedded runtime and board-level validation}

% T77-SOURCE: T75-FIG-A01; T76-PREVIEW-A01; T74-FIG-03; T74-TBL-05; T74-TBL-06; T74-SUP-03; T74-SUP-04; review=T72,T76
Embedded inference and board-level validation remain separate from software-HIL
simulation.  Current host evidence stops at two conservative facts: selected
preserved \texttt{.tflite} artifacts run in an isolated current-host
environment, and the read-only real-board gate/regeneration/provenance lane
still ends in
\texttt{NO\_GO\_REAL\_BOARD\_HOST\_OR\_DEVICE\_PATH\_UNAVAILABLE}.
Exported-model equivalence, embedded runtime latency, host-to-device update
cost, hardware resource use, and closed-loop timing therefore remain future
evidence rather than completed results.

\subsection{Parallel sidecar extension waves}

The main benchmark and paper-material lane should remain separated from
parallel extension experiments.  A useful way to organize the next experiments
is to treat them as sidecar waves: each wave may run in parallel with the
mainline, but its outputs remain sidecar evidence until a later promotion gate
decides whether a result deserves a new mainline task.  In particular, these
experiments must not rewrite the frozen four-scenario anchor, promote the
statistical-calibration lane into a mature comparator, or turn software-HIL
evidence into \texttt{.tflite} or board-level validation.

Wave A contains the immediately bounded sidecar lanes.  These lanes preserve
the fast-loop contract and mostly operate through cached replay, toy
simulation, or deterministic contract tests.

\begin{center}
\small
\begin{tabular}{>{\raggedright\arraybackslash}p{0.24\linewidth}
>{\raggedright\arraybackslash}p{0.39\linewidth}
>{\raggedright\arraybackslash}p{0.27\linewidth}}
\toprule
Wave A lane & Experimental question & First safe execution \\
\midrule
Temporal histogram TCN & Does a short causal sequence of histogram windows,
deltas, and compact statistics improve residual-\(b\) or drift-regime
prediction beyond the current short-window features? & Cached replay or tiny
TCN residual-head probe under a bounded sidecar run root. \\
Adaptive teacher confidence & Can a syndrome-only or teacher-output replay
estimate confidence, freeze unsafe updates, or fall back to a conservative
anchor before a bad parameter commit? & Offline teacher-policy replay with
explicit confidence and fallback fields. \\
Gain-scheduled parameter bank & Can a finite piecewise-affine bank reduce
parameter jitter or improve drift tracking while keeping the fast path as
\(\Delta=K_i s+b_i\)? & Bank-selection toy replay with range, thrashing, and
switch-frequency diagnostics. \\
Atomic commit and rollback & Which stage, commit, acknowledgement, freeze,
rollback, and timeout states are needed before stronger runtime integration is
credible? & Deterministic mock contract tests only; no MMIO, DMA, or board
writes. \\
\bottomrule
\end{tabular}
\end{center}

Wave B is the second-wave candidate pool.  These routes are useful for making
the paper experimentally complete, but they should start as S0 design or toy
probes rather than as long formal benchmarks.

\begin{center}
\small
\begin{tabular}{>{\raggedright\arraybackslash}p{0.25\linewidth}
>{\raggedright\arraybackslash}p{0.40\linewidth}
>{\raggedright\arraybackslash}p{0.25\linewidth}}
\toprule
Wave B lane & Experimental question & Boundary \\
\midrule
Moments and EWMA feature probes & Do low-cost summaries such as marginals,
mean/variance drift, EWMA, entropy, anisotropy, or edge-mass features explain
the same gains as the full histogram image? & Feature ablation only; not a new
formal benchmark by itself. \\
Teacher-conditioned residual head & Does conditioning the residual head on
teacher state, teacher deltas, or confidence improve stability without making
the learned module the per-shot decoder? & Small retrain or replay; still a
bounded slow-loop calibration module. \\
State-space or lightweight sequence model & Are modern compressed temporal
models such as tiny state-space, Mamba-like, or recurrent/transformer toy
heads useful after the TCN baseline is known? & Research-first toy route; no
full decoder replacement and no fast-loop change. \\
Outer-code soft-information positioning & How would the physical GKP
calibration layer interface with surface-GKP, GKP-surface, or QLDPC-GKP
soft-information decoders? & Positioning memo or tiny interface toy; not a
current mainline code-family benchmark. \\
\bottomrule
\end{tabular}
\end{center}

The purpose of this wave structure is not to multiply unrelated experiments.
It keeps the scientific question fixed: the fast path remains a deterministic
affine correction, while sidecar lanes ask which slow-loop estimator, feature
summary, teacher policy, or commit policy most safely updates the affine
parameters under drift.  Negative or inconclusive sidecar results are also
useful, because they can show that the current histogram-conditioned affine
contract is already near the useful complexity boundary.

% T80-REOPEN: Discussion
% T78-SCOPE: evidence-alignment; source=T77-manifest,T70,T72,T77
\section{Discussion}

The results and architecture suggest a useful reframing of the project.  The
central claim should not be that a CNN alone is a universally superior GKP
decoder.  A more defensible claim is that drift-adaptive GKP correction can be
organized as a low-dimensional affine calibration problem with a deterministic
fast path and replaceable slow calibration modules.  Under this framing, the
teacher-anchored CNN branch, the statistical calibration branch, and the
filtering baselines are all instances of the same runtime question: how should
recent syndrome statistics update \((K,b)\) without disturbing the per-shot
latency contract?

This framing clarifies why the method is different from several neighboring
approaches.  It is not a conventional FPGA surface-code decoder, because it
does not operate on a binary detection-event graph.  It is not an outer-code
GKP soft-information decoder, because it updates the physical displacement
rule before outer-code message passing or matching.  It is not an end-to-end
neural QEC decoder, because the neural component is not placed on the
per-shot critical path.  It is also not simply a calibration pre-processing
routine, because the estimated state is committed through the same quantized,
clipped, double-buffered runtime interface that the fast loop consumes.

% T77-SOURCE: T74-TBL-01; T74-TBL-07; T74-SUP-01; T75-FIG-M01; T75-FIG-M02; review=T70,T76
The advantage should be interpreted at two levels.  At the mainline project
level, the strongest accepted result is still the locked four-scenario
software-HIL ranking: the teacher-anchored hybrid branch beats the filtering
baselines across all frozen scenarios.  At the model-comparison level, the
current evidence does not require claiming that the CNN is always the best
estimator.  A bounded statistical-calibration extension lane is instead useful
because it shows that the same runtime contract can host a strong non-neural
alternative; however, the accepted closure still keeps that lane supplement-side
and no-promotion because the default and high-threshold settings remain
persistently tied.  The learned branch therefore remains valuable as a
replaceable slow-loop module rather than as a uniquely validated final
comparator.

The expected long-term advantage is therefore modularity under real-time
constraints.  If a future device requires a stronger statistical estimator, a
learned temporal model, or a more conservative fallback policy, the slow
calibration module can be replaced without changing the affine hardware
boundary.  Conversely, if future GKP experiments require tighter latency, the
fast path can remain a fixed-point matrix-vector operation while the slow
estimator is improved offline or on a host processor.  This is the main
systems-level reason to prefer the proposed contract over a monolithic decoder
design: the architecture can trade estimator sophistication against update
cadence without putting neural inference or a graph decoder on the per-shot
critical path.

% T77-SOURCE: T75-FIG-A01; T76-PREVIEW-A01; T74-TBL-05; T74-TBL-06; T74-SUP-03; T74-SUP-04; review=T72,T76
The boundary of the current evidence remains important.  Software-HIL results
support the adaptive-affine calibration story under controlled drift
scenarios.  Supporting materials show that the canonical training chain is
intact and that one clean CPU-only bounded rerun has been completed, but they
do not establish full reproducibility portability.  Isolated current-host
\texttt{.tflite} execution supports a narrow exported-model runtime claim, and
read-only real-board gate, regeneration, and provenance checks preserve a
current-host \texttt{NO\_GO} hardware boundary.  These layers do not yet
establish board-level latency, FPGA resource use, or closed-loop feedback
timing.  Those measurements are the natural next evidence layer if the work is
to move from ``hardware-compatible'' to hardware-demonstrated.

% T80-REOPEN: Conclusion
% T78-SCOPE: evidence-alignment; source=T77-manifest,T70,T72,T77
\section{Conclusion}

This work frames drift-adaptive GKP correction as a dual-loop affine
calibration problem.  The fast loop remains a bounded fixed-point rule,
\(\Delta_t=K_t s_t+b_t\), while a slower loop updates the affine parameters
from recent syndrome statistics.  This formulation creates a practical middle
ground between fixed decoders, which cannot track hardware drift, and full
neural decoders, which can be difficult to place on a deterministic per-shot
path.

% T77-SOURCE: T74-TBL-01; T74-TBL-02; T74-TBL-07; T74-SUP-01; review=T70,T76
The numerical results support two main conclusions.  First, under the locked
four-scenario software-HIL suite, adaptive affine decoding with a
teacher-anchored residual branch improves over filtering baselines, while
bounded ablations and multi-seed mechanism materials show that histogram-delta
features carry useful information and that teacher-side inputs do not admit a
simple necessity story.  Second, a separately labeled supplement-side
statistical-calibration extension lane shows that simple non-neural estimators
can also be strong under the same runtime contract, but the accepted closure
remains no-promotion and no unique clean threshold.  Together, these results
suggest that the strongest scientific story is not that one CNN architecture
has closed the comparison, but that low-dimensional histogram-driven
calibration of an affine GKP fast path is a promising structure for
drift-adaptive decoding.

% T77-SOURCE: T75-FIG-A01; T76-PREVIEW-A01; T74-TBL-05; T74-TBL-06; T74-SUP-03; T74-SUP-04; review=T72,T76
The remaining technical gap is to connect this bounded software-HIL evidence to
broader drift families, stronger oracle and wrapped-Gaussian baselines,
measured runtime constraints, and eventually real hardware timing.  Current
supporting and deployment-facing materials remain layered: training/material
evidence reaches a canonical chain plus one clean CPU-only bounded rerun,
isolated true runtime is verified only on the current host, and the real-board
lane remains gate-only and \texttt{NO\_GO}.  Future work must therefore test
whether the affine calibration contract remains useful when the benchmark moves
from a controlled drift suite toward a deployment-grade decoder evaluation.  If
successful, the resulting contribution would be a practical bridge between
analog GKP decoding theory and low-latency FPGA-QEC systems: a physical-layer
calibration module that preserves analog information, tracks drift, and still
exposes a compact deterministic fast path.

\begin{thebibliography}{99}

\bibitem{gkp2001}
D. Gottesman, A. Kitaev, and J. Preskill,
``Encoding a qubit in an oscillator,''
\emph{Physical Review A}, 2001.

\bibitem{grimsmo2021}
A. L. Grimsmo and S. Puri,
``Quantum Error Correction with the Gottesman--Kitaev--Preskill Code,''
\emph{PRX Quantum} 2, 020101, 2021.

\bibitem{fukui2018}
K. Fukui, A. Tomita, A. Okamoto, and K. Fujii,
``High-threshold fault-tolerant quantum computation with analog quantum error
correction,''
\emph{Physical Review X}, 2018.

\bibitem{noh2020}
K. Noh and C. Chamberland,
``Fault-tolerant bosonic quantum error correction with the surface-GKP code,''
\emph{Physical Review A}, 2020.

\bibitem{noh2022}
K. Noh, C. Chamberland, and F. G. S. L. Brandao,
``Low-overhead fault-tolerant quantum error correction with the surface-GKP
code,''
\emph{PRX Quantum}, 2022.

\bibitem{raveendran2022}
N. Raveendran, N. Rengaswamy, F. Rozpedek, A. Raina, L. Jiang, and B. Vasic,
``Finite-rate QLDPC-GKP coding scheme that surpasses the CSS Hamming bound,''
\emph{Quantum}, 2022.

\bibitem{berent2024}
L. Berent et al.,
``Analog information decoding of bosonic quantum LDPC codes,''
\emph{PRX Quantum}, 2024.

\bibitem{borah2025}
S. K. Borah, A. K. Pradhan, N. Raveendran, M. Pacenti, and B. Vasic,
``Fault tolerant decoding of QLDPC-GKP codes with circuit level soft
information,''
preprint, 2025.

\bibitem{roy2025}
M.-A. Roy, T. Pousset, and B. Royer,
``Decoding multimode Gottesman--Kitaev--Preskill codes with noisy auxiliary
states,''
preprint, 2025.

\bibitem{spitz2018}
S. Spitz, B. Tarasinski, C. W. J. Beenakker, and T. E. O'Brien,
``Adaptive weight estimator for quantum error correction,''
\emph{Advanced Quantum Technologies}, 2018.

\bibitem{wagner2021}
T. Wagner, H. Kampermann, and D. Bru{\ss},
``Optimal noise estimation from syndrome statistics of quantum codes,''
\emph{Physical Review Research}, 2021.

\bibitem{dgr2023}
H. Chen et al.,
``DGR: Tackling drifted and correlated noise in quantum error correction via
decoding graph re-weighting,''
preprint, 2023.

\bibitem{chen2022}
E. H. Chen et al.,
``Calibrated decoders for experimental quantum error correction,''
\emph{Physical Review Letters}, 2022.

\bibitem{sivak2024}
V. Sivak, M. Newman, and P. Klimov,
``Optimization of decoder priors for accurate quantum error correction,''
\emph{Physical Review Letters} 133, 150603, 2024.

\bibitem{chamberland2026}
C. Chamberland, J. Olle, M. Li, S. Thornton, and I. Baratta,
``Fast and accurate AI-based pre-decoders for surface codes,''
arXiv:2604.12841, 2026.

\bibitem{stein2026}
S. Stein, S. Kan, C. Liu, A. Harkness, S. Garner, Z. Du, Y. Ding, Y. Mao, and
A. Li,
``Calibration-conditioned FiLM decoders for low-latency decoding of quantum
error correction evaluated on IBM repetition-code experiments,''
preprint, 2026.

\bibitem{google2023decoder}
Google Quantum AI and collaborators,
``Learning high-accuracy error decoding for quantum processors,''
preprint, 2023.

\bibitem{lilliput2022}
``LILLIPUT: A lightweight low-latency lookup-table based decoder for near-term
quantum error correction,''
\emph{ASPLOS}, 2022.

\bibitem{afs2022}
``AFS: Accurate, fast, and scalable error-decoding for fault-tolerant quantum
computers,''
\emph{HPCA}, 2022.

\bibitem{helios2023}
``Scalable quantum error correction for surface codes using FPGA,''
preprint, 2023.

\bibitem{astrea2023}
``Astrea: Accurate quantum error-decoding via practical minimum-weight
perfect-matching acceleration,''
\emph{ISCA}, 2023.

\bibitem{collision2025}
``A real-time, scalable, fast and resource-efficient decoder for a quantum
computer,''
\emph{Nature Electronics}, 2025.

\bibitem{qldpcfpga2025}
``FPGA-tailored algorithms for real-time decoding of quantum LDPC codes,''
preprint, 2025.

\bibitem{caune2024}
L. Caune et al.,
``Demonstrating real-time and low-latency quantum error correction with
superconducting qubits,''
preprint, 2024.

\bibitem{yang2026}
X. Yang et al.,
``Real-time surface-code error correction using an FPGA-based neural-network
decoder,''
preprint, 2026.

\bibitem{ziad2024}
M. Ziad et al.,
``Local clustering decoder: a fast and adaptive hardware decoder for the
surface code,''
preprint, 2024.

\bibitem{maurer2025}
L. Maurer et al.,
``Real-time decoding of the gross code memory with FPGAs,''
preprint, 2025.

\bibitem{sivak2023}
V. Sivak et al.,
``Real-time quantum error correction beyond break-even,''
\emph{Nature}, 2023.

\bibitem{lachance2024}
D. Lachance-Quirion et al.,
``Autonomous quantum error correction of Gottesman--Kitaev--Preskill states,''
\emph{Physical Review Letters}, 2024.

\end{thebibliography}

\end{document}
